\begingroup\fontsize{8}{10}\selectfont

\begin{longtable}[t]{lp{4.2cm}p{8.2cm}}
\caption{\label{tab:appD_quotes_part1}Open-ended responses: methodology and marginal cost. Verbatim quotes; empty entries and pure `N/A' answers are omitted.}\\
\toprule
Section & Prompt & Response\\
\midrule
\endfirsthead
\caption[]{Open-ended responses: methodology and marginal cost. Verbatim quotes; empty entries and pure `N/A' answers are omitted. \textit{(continued)}}\\
\toprule
Section & Prompt & Response\\
\midrule
\endhead

\endfoot
\bottomrule
\endlastfoot
Cost curvature & Additional robustness checks or alternative specifications & Not sure\\
 &  & I need to spend more time with the paper to think about this.\\
 &  & Wind speed duration is a significant factor as well, and can be computed from the ensemble wind fields (but not from the swath, you'd have to get the time-evolving wind fields).\\
 &  & Hurricane force wind gusts can occur far outside of the sustained hurricane wind areas and can do significant damage far inland. The 2024 landfall of Hurricane Helene was a recent example, where wind damage occurred far inland to the Carolinas and even Virginia, but no sustained hurricane winds were observed beyond central Georgia.\\
 &  & Slides presentations are good for an introduction but reviewing the methodology in writing is necessary for the level of questions you are asking. My neutral response could be higher with more time to read and think about it.\\
\addlinespace
 &  & Track, precip, and storm surge should be accounted for. You can also test the storms minimum central pressure rather than maximum winds: some papers suggest this metric is more robust. Weak TCs can also cause extensive flood damages.\\
 &  & unsure\\
 &  & None.\\
Cost curvature & Conceptual or empirical gaps in how forecast error is monetised & No\\
 &  & I need to spend more time with the paper to think about this.\\
\addlinespace
 &  & No.\\
 &  & This is very complicated problem and I don't have enough experience in this area to comment.\\
 &  & It may not describe the full sample space of benefits.\\
 &  & No\\
 &  & I honestly don't understand the strategy - not a limitation of the method, I don't think, but a limitation of my own understanding.\\
\addlinespace
 &  & Damages are based on observed hurricane impacts with status quo forecasts but we don't have a counterfactual of damages with improved forecasts. Its unclear whether the reductions in damages from forecast improvement would be realized by on-the-ground changes in behavior (protective actions, better preparation/sheltering). I don't believe this is a causal relationship, just a statistical finding. So the interpretation, I believe, is that the paper provides a value of forecast improvement but it shouldn't be interpreted as a reduction in damages attributable to forecast improvement.\\
Methodology & Additional dimensions of forecast quality that should be included & Surge would be nice to have, but understand the challenges.\\
 &  & Storm surge\\
 &  & Absolutely, all the above. Wind speed is not the killer of public, it's water. Expansion on track as well as water hazards, rainfall and surge should be considered.\\
 &  & Precipitation is important and harder to forecast than wind speed.\\
\addlinespace
 &  & Track is fundamental to wind and surge hazards so it would be included in the future. Storm surge should also be added. Precipitation is trickier because there are many factors that impact precip uncertainty, including those not directly related to the hurricane such as PREs and other interactions with baroclinic weather systems.\\
 &  & Track, surge, precipitation are also considered by the NOAA Chief Economist where value chain linkages can be confirmed (aka SLOSH model saves lives).\\
 &  & The three metrics listed above should all be included! Andrew Martinez has done some work on using TC track errors to estimate damages. His work can be a good starting place. Storm surge and freshwater flooding have been shown to cause high damage amounts. Untangling the impacts of these features could be challenging. Flooding from weak TCs (Harvey, ida, Helene, etc) especially  for inland regions should be considered.\\
 &  & Yes, all of the above, ideally.\\
 &  & Track and storm surge could be assessed in the same fashion to determine significance in cost reduction of hurricane damages.\\
\addlinespace
Methodology & Methodological limitations linking augmented observations to forecast precision & No\\
 &  & I am not as familiar as to comment further on methodological limitations at this time.\\
 &  & No.\\
 &  & Because hurricanes are relatively rare events, especially for specific locations, sample size will continue to be a limitation.\\
 &  & Data-to-decision pathways are very hard to lock down. Many NOAA professionals do not understand the full breadth of how observational data transforms into models and then is derived into a decisional product which is then used by a societal sector.\\
\addlinespace
 &  & No, I think you're on the right track!\\
 &  & No\\
 &  & It seems reasonable.\\
Methodology & On whether the 0.26 m/s reduction is operationally meaningful & It seems like a small value, but averaging over storms and counties would do that.\\
 &  & I don't think I have enough information to comment further on this question.\\
\addlinespace
 &  & Any improvements is valuable even less than 1 m/s, especially if you are the individual impacted by the storm, I.e. insurance claims, life, property. Numbers matter.\\
 &  & Sounds kind of small, but when averaged over many counties and many storms, I think it will easily be statistically significant.\\
 &  & The response to that change seems very large, so even if the qualitative results have large error bars, the qualitative results are still meaningful.\\
 &  & You should provide a "I don't know" category for those not subject matter experts. You don't ask this, but 3 days seems to be the lower bound for significant tropical storm/hurricane decisionmaking. I understand some decisions occur further out.\\
 &  & I expected that dropsonde observations would've had a bigger impact on the wind reduction error during landfalling TCs. While still meaningful, it seems more challenging to justify such a small number as the net impact of our observations.\\
\addlinespace
 &  & I don't have a good feel for the meaning of this metric.  It seems like a small number on its face, but that's in the context of predicting maximum sustained winds. I think to understand this number better, I need to understand the improvement in the forecast over the 15-year period.\\
 &  & I'm not qualified to answer this.\\*
\end{longtable}
\endgroup{}
